\documentclass[12pt,a4paper]{article}
\usepackage[utf8]{inputenc}
\usepackage[spanish,es-noquoting]{babel}
\usepackage[T1]{fontenc}
\usepackage{lmodern}
\usepackage{geometry}
\geometry{top=2cm, bottom=2cm, left=2.5cm, right=2.5cm}
\usepackage{xcolor}
\usepackage{listings}
\usepackage{enumitem}
\usepackage{fancyhdr}
\usepackage{hyperref}

\definecolor{codegreen}{rgb}{0,0.6,0}
\definecolor{backcolour}{rgb}{0.95,0.95,0.92}

\lstset{
    language=C,
    backgroundcolor=\color{backcolour},
    commentstyle=\color{codegreen},
    keywordstyle=\color{blue},
    basicstyle=\ttfamily\small,
    breaklines=true,
    frame=single,
    numbers=left,
    tabsize=4,
    extendedchars=true,
    showstringspaces=false,
    literate={á}{{\'a}}1 {é}{{\'e}}1 {í}{{\'i}}1 {ó}{{\'o}}1 {ú}{{\'u}}1 {ñ}{{\~n}}1
}

\pagestyle{fancy}
\fancyhf{}
\lhead{Monitorías Sistemas Operativos 2026-I}
\rhead{Capítulo 0: Ejercicios de Repaso}
\cfoot{\thepage}

\title{\textbf{Taller Práctico: Memoria y Archivos}}
\author{Malak Sanchez \\ \small Monitorías de Sistemas Operativos}
\date{}

\begin{document}

\maketitle

\section*{Introducción}
El siguiente taller tiene como objetivo integrar los conceptos fundamentales vistos en el Capítulo 0: paso de parámetros por línea de comandos, gestión de memoria dinámica y manipulación de archivos. Los problemas planteados requieren un análisis detallado de la estructura de datos y una gestión responsable de los recursos del sistema.

\section*{\mbox{Problema 1: Generador Dinámico de Patrones de Marco}}
Diseñe un programa que genere una matriz ``marco'' en memoria dinámica. El programa debe recibir por línea de comandos dos enteros positivos $N$ y $M$, y un carácter $C$.

\textbf{Ejemplo de ejecución:} 
\begin{lstlisting}[numbers=none]
$ ./pattern_gen 5 5 #
\end{lstlisting}

\textbf{Ejemplo de salida:}
\begin{lstlisting}[numbers=none]
# # # # #
# . . . #
# . . . #
# . . . #
# # # # #
\end{lstlisting}

\textbf{Requerimientos:}
\begin{enumerate}
    \item El programa debe validar que se pasen los 3 argumentos necesarios.
    \item Debe reservar memoria dinámica para una matriz de caracteres (\texttt{char **}) de tamaño $N \times M$.
    \item Rellene los bordes de la matriz con el carácter $C$ y el interior con puntos (\texttt{.}).
    \item Implemente una función \texttt{print\_matrix} y una función \texttt{free\_matrix} para gestionar la limpieza.
    \item El programa debe imprimir la matriz en pantalla.
    \item La matriz creada debe ser dinámica, es decir, debe poder aceptar cualquier valor de $N$ y $M$ e imprimir correctamente el patrón.
\end{enumerate}

\newpage

\section*{\mbox{Problema 2: El Analizador de Bitácoras (Log Parser)}}
Diseñe una herramienta que procese un archivo de texto con registros de eventos de un sistema. La primera fila del archivo contien un valor $n$ que nos dice la cantidad de elementos a leer. Cada línea del archivo representa un evento con el formato:

\begin{lstlisting}[numbers=none]
<ID_PROCESO> <TIMESTAMP> <CONSUMO_MEMORIA_MB>
\end{lstlisting}

\textbf{Ejemplo de comando:} 
\begin{lstlisting}[numbers=none]
$ ./log_parser input.txt 100
\end{lstlisting}

\textbf{Ejemplo de salida esperada (Consumo superior a 100 MB):} 
\begin{lstlisting}[numbers=none]
1024 1710432000 101.5
1025 1710432001 120.0
1026 1710432002 200.0
\end{lstlisting}

\textbf{Requerimientos:}
\begin{enumerate}
    \item Debe leer el archivo y almacenar los datos correctamente.
    \item Debe usar \texttt{fscanf(file, "\%d \%f \%f", \&id, \&timestamp, \&memory\_consumption)} para leer los datos del archivo con el formato especificado.
    \item Filtre y muestre solo los procesos cuyo consumo de memoria sea superior al umbral pasado por parámetro.
    \item Cuente la cantidad de procesos que consumen más memoria que el umbral e imprima el resultado. 
\end{enumerate}

\newpage

\section*{\mbox{Problema 3: Mini-Compresor de Texto RLE}}
Implemente un programa que reciba como parametros un archivo $input.in$ y realice una compresión por longitud de caracteres (RLE - Run-Length Encoding).

\textbf{Ejemplo de comando:}
\begin{lstlisting}[numbers=none]
$ ./rle_compress input.in
\end{lstlisting}

\textbf{Ejemplo visual de funcionamiento:}
\begin{lstlisting}[numbers=none]
Entrada (en archivo): AAAAABBBCC
Salida (consola): 5A3B2C
\end{lstlisting}

\textbf{Requerimientos:}
\begin{enumerate}
    \item Cargue el contenido del archivo pasado como argumentos del main.
    \item El programa debe validar que el archivo de entrada no esté vacío antes de proceder.
    \item Almacene la cadena comprimida en un \textit{buffer}.
    \item Imprima en pantalla el ``Ratio de Compresión'' (Tamaño Original / Tamaño Comprimido).
    \item Libere correctamente la memoria dinámica utilizada.
\end{enumerate}

\section*{Criterios de Evaluación}
\begin{itemize}
    \item \textbf{Gestión de Memoria:} Uso correcto de \texttt{malloc/free} y ausencia de fugas.
    \item \textbf{Robustez:} El programa no debe fallar si el archivo no existe o los argumentos son insuficientes.
    \item \textbf{Claridad:} Código bien estructurado y documentado en lugares clave.
\end{itemize}

\end{document}
